The careful work with literature is the basis for every publication in all
branches of science. This is not just a matter of compiling a bibliography. The
real work begins even before you start writing your own publication. The
following steps are necessary for the creation of bibliographies.
%
\begin{compactenum}
\item Literature research and creation of one or several litature databases.
      This step is independent of the LaTeX workflow.
\item Integration of the literature databases into the \LaTeX document.
\item Creation of literature references in the \LaTeX document.
\item Adaptation of the design of the bibliography, if necessary.
\end{compactenum}

\LaTeX\ offers many ways to create a bibliography. In the simplest
case, the bibliography is created manually by the \texttt{thebibliography} environment.
This approach leads to quick results and is quite practicable for smaller
seminar papers. For more extensive work, however, the effort quickly gets out
of hand. In the context of this project we therefore use the additional
package~\texttt{biblatex}, which offers the following advantages.
%
\begin{compactitem}
\item Once a literature database has been created, it can be reused in later
      publications.
\item Literature markers are automatically created using predefined schemes.
\item The formatting of the bibliography is carried out uniformly according to
      relevant specifications.
\item Only those bibliographical references from the database are included in the
      The bibliography is always formatted in accordance with the relevant specifications.
\end{compactitem}

The design of bibliographies is a complex subject. In the context of this
instructions, we will only discuss how the bibliography can be integrated into
the template technicaly. For detailed questions or more information on the
layout, please refer to \cite{voss_bibliografien} and \cite[chapter~13]{voss_wissenschaftliche_arbeit}.
The package~\ctanref{biblatex} has a wide range of settings, for which we refer
to the manual. A brief overview of the package and the structure of the
literature database can be found under \cite{rees_biblatex_cheat}.
